\documentclass[aspectratio=169,10pt]{beamer}
% Shared Beamer setup for Fall 2026 course slides.
% Clean Cal Poly-inspired theme: no custom class files or uncommon packages.
\mode<presentation>{
  \usetheme{default}
  \usefonttheme{professionalfonts}
}

\definecolor{CalPolyGreen}{HTML}{154734}
\definecolor{CalPolyGold}{HTML}{BD8B13}
\definecolor{CalPolySage}{HTML}{7FA28F}
\definecolor{CalPolyMint}{HTML}{DDF1E7}
\definecolor{CalPolySand}{HTML}{A29F86}
\definecolor{CalPolyBeige}{HTML}{EFEEDE}
\definecolor{DeckInk}{HTML}{1F2933}
\definecolor{DeckMuted}{HTML}{5F6B7A}
\definecolor{DeckSoft}{HTML}{F7F7F5}

\setbeamersize{text margin left=0.55in,text margin right=0.55in}
\setbeamercolor{normal text}{fg=DeckInk,bg=white}
\setbeamercolor{structure}{fg=CalPolyGreen}
\setbeamercolor{title}{fg=white}
\setbeamercolor{frametitle}{fg=CalPolyGreen,bg=white}
\setbeamercolor{block title}{bg=CalPolySage,fg=white}
\setbeamercolor{block body}{bg=CalPolyMint,fg=black}
\setbeamercolor{alerted text}{fg=CalPolyGold}

\setbeamerfont{title}{size=\Huge,series=\bfseries}
\setbeamerfont{frametitle}{size=\Large,series=\bfseries}
\setbeamerfont{block title}{size=\normalsize,series=\bfseries}

\setbeamertemplate{navigation symbols}{}
\setbeamertemplate{itemize item}{\color{CalPolyGold}$\blacktriangleright$}
\setbeamertemplate{itemize subitem}{\color{CalPolyGreen}$\bullet$}
\setbeamertemplate{footline}{
  \hfill{\color{DeckMuted}\scriptsize\insertframenumber/\inserttotalframenumber}
  \hspace{0.18in}\vspace{0.12in}
}
\setbeamertemplate{frametitle}{
  \vspace{0.18in}
  \hspace{0.02in}\insertframetitle\par
  \vspace{0.08in}
  \noindent\makebox[\linewidth][l]{%
    {\color{CalPolyGold}\rule{0.55in}{2pt}}%
    {\color{CalPolyGreen}\rule{\dimexpr\linewidth-0.55in\relax}{2pt}}%
  }
}

\newcommand{\CourseNumber}{Course}
\newcommand{\CourseTitle}{Course Title}
\newcommand{\LectureNumber}{0}
\newcommand{\LectureTitle}{Lecture Title}
\newcommand{\LectureDate}{Date}
\newcommand{\CourseUnit}{Unit}

\newcommand{\coursenumber}[1]{\renewcommand{\CourseNumber}{#1}}
\newcommand{\coursetitle}[1]{\renewcommand{\CourseTitle}{#1}}
\newcommand{\lecturenumber}[1]{\renewcommand{\LectureNumber}{#1}}
\newcommand{\lecturetitle}[1]{\renewcommand{\LectureTitle}{#1}}
\newcommand{\lecturedate}[1]{\renewcommand{\LectureDate}{#1}}
\newcommand{\courseunit}[1]{\renewcommand{\CourseUnit}{#1}}

\author{Kirk Duran}
\institute{California Polytechnic State University, San Luis Obispo}
\date{\LectureDate}

\newcommand{\makedecktitle}{
  \begingroup
  \setbeamercolor{background canvas}{bg=CalPolyGreen}
  \begin{frame}[plain]
    \vfill
    \begin{center}
      {\usebeamerfont{title}\color{white}\LectureTitle\par}
      \vspace{0.18in}
      {\Large\color{white}Lecture \LectureNumber\par}
    \end{center}
    \vfill
    \noindent{\color{white}\rule{\linewidth}{1.2pt}}\par
    \vspace{0.08in}
    \noindent\colorbox{CalPolyGold}{\makebox[0.24\linewidth][c]{\color{white}\Large\CourseNumber}}
    \hspace{0.12in}{\color{white}\Large Kirk Duran}
  \end{frame}
  \endgroup
}

\newenvironment{keyidea}{\begin{block}{Key idea}}{\end{block}}
\newenvironment{activity}{\begin{block}{In class}}{\end{block}}
\newenvironment{cpdefinition}{\begin{block}{Definition}}{\end{block}}
\newenvironment{exampleblockcp}{
  \begingroup
  \setbeamercolor{block title}{bg=CalPolySand,fg=white}
  \setbeamercolor{block body}{bg=CalPolyBeige,fg=black}
  \begin{block}{Example}
}{
  \end{block}
  \endgroup
}

\newcommand{\imageplaceholder}[2][0.3in]{%
  \noindent\makebox[\linewidth][c]{%
    \fcolorbox{CalPolySage}{CalPolyBeige}{%
      \parbox[c][#1][c]{0.86\linewidth}{%
        \centering
        {\color{DeckMuted}\scriptsize Image placeholder: }%
        {\color{DeckInk}\scriptsize\ttfamily\detokenize{#2}}%
      }%
    }%
  }%
}

\newcommand{\sectionframe}[1]{
  \begingroup
  \setbeamercolor{background canvas}{bg=CalPolyGreen}
  \begin{frame}[plain]
    \vfill
    {\color{white}\LARGE\bfseries #1\par}
    \vspace{0.18in}
    {\color{CalPolyGold}\rule{0.22\paperwidth}{3pt}}
    {\color{white}\rule{0.62\paperwidth}{3pt}}
    \vfill
  \end{frame}
  \endgroup
}

\newcommand{\placeholdernote}{
  \begin{block}{Placeholder}
    Replace these bullets with the final lecture material, examples, and in-class activities.
  \end{block}
}

\usepackage[T1]{fontenc}
\usepackage{lmodern}
\usepackage{amsmath,amssymb}
\usepackage{booktabs,array}
\pdfstringdefDisableCommands{\def\translate#1{#1}}
\graphicspath{{../assets/lecture-29/}}
% Adapted from instructor-supplied Lecture 20 (40-page) neural-networks PDF.
% Lecture sequence continues Lectures 26-28. Preserve instructor examples:
% breast-cancer perceptron, Keras breast-cancer MLP, sine regression, MNIST.
% Source historical story preserved with careful qualifications.
% Approx. 50 minutes: 5 motivation, 10 perceptrons/XOR, 10 MLP/activation,
% 9 backprop/history, 8 cancer practical, 6 sine/MNIST, 2 limitations/synthesis.
% Technical support in calculation boxes; students need not know advanced calculus.
% Checked against Keras Model training APIs, scikit-learn load_breast_cancer,
% and Rumelhart, Hinton & Williams (1986), https://www.cs.toronto.edu/~hinton/absps/naturebp.pdf
% https://keras.io/api/models/model_training_apis/
% https://scikit-learn.org/stable/modules/generated/sklearn.datasets.load_breast_cancer.html
\newcommand{\lecturevisual}[3][2.15in]{%
 \par\vspace{2pt}{\centering\includegraphics[width=.98\linewidth,height=#1,keepaspectratio]{#2}\par}\vspace{2pt}%
}
\newcommand{\conceptlead}[1]{{\small\bfseries\color{CalPolyGreen}#1\par}}
\newenvironment{calculation}[1][Mathematical description]{%
 \par\begingroup\setbeamercolor{calculation body}{bg=DeckSoft,fg=DeckInk}%
 \begin{beamercolorbox}[sep=3pt,wd=\linewidth]{calculation body}%
 \footnotesize{\bfseries\color{DeckMuted}#1\par}%
 \setlength{\abovedisplayskip}{3pt}\setlength{\belowdisplayskip}{3pt}\setlength{\jot}{0pt}%
}{\end{beamercolorbox}\endgroup}
\setbeamersize{text margin left=.42in,text margin right=.42in}
\coursenumber{CS 4880}
\coursetitle{Artificial Intelligence}
\lecturenumber{29}
\lecturetitle{Neural Networks: From Perceptrons to Deep Learning}
\lecturedate{November 4, 2026}
\courseunit{Learning Agents and Machine Learning}
\begin{document}
\makedecktitle

\section{Why Another Model Family?}
\begin{frame}[t]{Where a Decision Tree Struggles}
\small
\conceptlead{Trees partition the world into regions; some targets change smoothly across those regions.}
\begin{columns}[T,onlytextwidth]
 \begin{column}{.48\textwidth}
 \begin{itemize}
  \item A standard decision tree predicts one value throughout each regression leaf: a \textbf{step-shaped approximation}.
  \item A curved or diagonal boundary may require many axis-aligned splits.
  \item Small changes in training records can produce a very different tree; ensembles reduce some of this instability.
 \end{itemize}
 \end{column}
 \begin{column}{.49\textwidth}
 \lecturevisual[1.95in]{tree_vs_smooth.png}{A depth-limited tree versus a smooth underlying function.}
 \end{column}
\end{columns}
\begin{keyidea}Neural networks learn connected numerical transformations rather than a sequence of hard feature tests.\end{keyidea}
\end{frame}

\begin{frame}[t]{Learning a Category Without a Rulebook}
\small
\conceptlead{We can recognize many different dogs without writing a rule for every breed.}
\begin{itemize}
 \item A child sees large, small, fluffy, and short-haired dogs; examples establish a useful category despite differences.
 \item Correlated clues such as shape, ears, tails, and movement can collectively support recognition.
 \item In supervised learning, we give a model examples with known labels; it learns adjustable combinations of clues.
 \item The model may learn statistical associations rather than the real-world meaning of ``dog.''
\end{itemize}
\begin{block}{Connect to learning agents}
The agent's \textbf{learning element} updates a predictor using feedback. Its \textbf{performance element} can then use that predictor to select actions. A neural network is a model inside this larger loop, not an entire learning agent by itself.
\end{block}
\end{frame}

\begin{frame}[t]{Biological Inspiration: Neurons and Plasticity}
\small
\conceptlead{The biological analogy is useful for intuition, but an artificial neuron is a mathematical simplification.}
\begin{columns}[T,onlytextwidth]
 \begin{column}{.54\textwidth}
 \begin{itemize}
  \item Biological neurons receive signals through connections, integrate inputs, and can produce an output signal.
  \item The strength of connections can change through experience, a phenomenon called \textbf{synaptic plasticity}.
  \item The Hebbian slogan ``neurons that fire together wire together'' captures one form of association learning.
 \end{itemize}
 \end{column}
 \begin{column}{.43\textwidth}
 \begin{block}{Artificial analogue}
 Inputs $\longrightarrow$ weighted sum $\longrightarrow$ activation $\longrightarrow$ output.
 \par\medskip
 Training modifies numerical \textbf{weights} to improve predictions. Modern backpropagation is not intended as a literal model of brain learning.
 \end{block}
 \end{column}
\end{columns}
\end{frame}

\section{Perceptrons: Learning a Linear Classifier}
\begin{frame}[t]{The Perceptron: A Trainable Artificial Neuron}
\small
\conceptlead{A perceptron makes a binary decision from a weighted vote of its inputs.}
\begin{columns}[T,onlytextwidth]
 \begin{column}{.53\textwidth}
 \lecturevisual[1.06in]{perceptron_structure.png}{Inputs, weights, bias, and threshold.}
 \end{column}
 \begin{column}{.43\textwidth}
 \begin{itemize}
  \item Feature weights increase or decrease influence.
  \item Add the weighted inputs and a bias (the decision's starting offset).
  \item Threshold the score to obtain 0 or 1.
 \end{itemize}
 \end{column}
\end{columns}
\begin{calculation}[Weighted score followed by a threshold]
\[z=b+\sum_{j=1}^{d}w_jx_j,\qquad \widehat y=\mathrm{step}(z).\]
$d$: feature count; $j$: feature index; $x_j$: input; $w_j$: weight; $b$: bias; $z$: weighted score; $\widehat y$: predicted label. Here $\mathrm{step}(z)=1$ when $z\geq0$, otherwise $0$.
\end{calculation}
\end{frame}

\begin{frame}[t]{How a Perceptron Learns from Mistakes}
\small
\conceptlead{When the label is wrong, adjust the weights toward the correct answer.}
\begin{enumerate}
 \item Compute the current prediction from the feature values.
 \item Compare the predicted label with the actual training label.
 \item If the prediction is correct, leave the weights unchanged for that example.
 \item If it is wrong, change each weight in proportion to its input feature.
\end{enumerate}
\begin{calculation}[Per-example perceptron update]
\[w_j\leftarrow w_j+\eta(y-\widehat y)x_j,\qquad
b\leftarrow b+\eta(y-\widehat y).\]
$y$: true 0/1 label; $\widehat y$: prediction; $\eta>0$: learning rate; $x_j$: feature; $w_j,b$: weight and bias. $\leftarrow$: replace with the updated value. The error $y-\widehat y$ is $+1$ for a missed class-1 example, $-1$ for a false class-1 prediction, and $0$ if correct.
\end{calculation}
\begin{keyidea}This is supervised learning: observed labels change the model rather than being stored as independent handwritten rules.\end{keyidea}
\end{frame}

\begin{frame}[fragile,t]{Prepare One Shared Train, Validation, and Test Split}
\small
\conceptlead{Use the same 30 features and partitions for the perceptron and the later MLP.}
\begin{columns}[T,onlytextwidth]
\begin{column}{.56\textwidth}
\scriptsize
\begin{verbatim}
from sklearn.datasets import load_breast_cancer
from sklearn.model_selection import train_test_split
from sklearn.preprocessing import StandardScaler
X, y = load_breast_cancer(return_X_y=True)
X_dev, X_test, y_dev, y_test = train_test_split(
    X, y, test_size=.2, stratify=y, random_state=42)
X_train, X_val, y_train, y_val = train_test_split(
    X_dev, y_dev, test_size=.2, stratify=y_dev,
    random_state=42)
\end{verbatim}
\end{column}
\begin{column}{.41\textwidth}
\scriptsize
\begin{verbatim}
scaler = StandardScaler().fit(X_train)
X_train = scaler.transform(X_train)
X_val = scaler.transform(X_val)
X_test = scaler.transform(X_test)
\end{verbatim}
\end{column}
\end{columns}
\begin{keyidea}Fit the scaler only on training rows. Dataset labels: $0$ = malignant; $1$ = benign. A sigmoid score for class 1 is a benign score, not a cancer-risk score.\end{keyidea}
\end{frame}

\begin{frame}[fragile,t]{Python Example: Breast-Cancer Perceptron}
\small
\conceptlead{Fit a linear classifier, predict held-out labels, then measure accuracy.}
\begin{columns}[T,onlytextwidth]
 \begin{column}{.52\textwidth}
{\footnotesize\begin{verbatim}
from sklearn.linear_model import Perceptron
from sklearn.metrics import accuracy_score

clf = Perceptron(max_iter=1000,
                 tol=1e-3,
                 random_state=42)
clf.fit(X_train, y_train)
y_pred = clf.predict(X_val)
acc = accuracy_score(y_val, y_pred)
print(f"Validation Accuracy: {acc:.2f}")
\end{verbatim}}
 \end{column}
 \begin{column}{.44\textwidth}
 \begin{block}{What each call means}
 \texttt{fit}: learn the feature weights.\\[5pt]
 \texttt{predict}: apply the learned boundary to held-out records.\\[5pt]
 \texttt{accuracy\_score}: count correct predictions.
 \end{block}
 \end{column}
\end{columns}
\begin{keyidea}The previous slide supplies the scaled data. Compare this baseline on validation; keep the final test out of model and epoch selection.\end{keyidea}
\end{frame}

\begin{frame}[t]{What Can That Perceptron Represent?}
\small
\conceptlead{The learned weights tilt and move one linear boundary through the feature space.}
\begin{columns}[T,onlytextwidth]
 \begin{column}{.54\textwidth}
 \lecturevisual[2.53in]{perceptron_cancer_boundary.png}{Example boundary using two breast-cancer features; not the 30-feature model.}
 \end{column}
 \begin{column}{.42\textwidth}
 \begin{itemize}
  \item This supplied plot uses mean radius and mean texture; it is not the 30-feature model fitted by our code.
  \item A single boundary cannot wrap around disconnected or highly curved clusters.
  \item Better features can help; adding hidden layers changes what the model can represent.
 \end{itemize}
 \end{column}
\end{columns}
\end{frame}

\begin{frame}[t]{XOR: A Simple Problem One Perceptron Cannot Solve}
\small
\conceptlead{Some patterns have no single straight line separating the labels.}
\begin{columns}[T,onlytextwidth]
 \begin{column}{.53\textwidth}
 \lecturevisual[2.35in]{xor_points.png}{The XOR positive class occupies opposite corners.}
 \end{column}
 \begin{column}{.43\textwidth}
 \begin{itemize}
  \item In \textbf{XOR}, exactly one input must be 1.
  \item The positive cases are $(0,1)$ and $(1,0)$; the negative cases are $(0,0)$ and $(1,1)$.
  \item Every proposed straight line misclassifies at least one corner.
 \end{itemize}
 \end{column}
\end{columns}
\begin{calculation}[XOR truth table]
$0\oplus0=0$, $0\oplus1=1$, $1\oplus0=1$, $1\oplus1=0$. $\oplus$ denotes exclusive OR; $x_1,x_2$ are binary inputs, $w_1,w_2$ are weights, and $b$ is bias. No single affine threshold matches all four cases.
\end{calculation}
\end{frame}

\begin{frame}[t]{The First Wave: A Breakthrough and a Boundary}
\small
\conceptlead{The perceptron made learning from data tangible, but its expressive limits were real.}
\begin{itemize}
 \item \textbf{1958:} Frank Rosenblatt's perceptron popularized trainable pattern recognition, building on earlier threshold-neuron models.
 \item \textbf{1969:} Minsky and Papert analyzed important limits of single-layer perceptron architectures, including XOR-like problems.
 \item Enthusiasm for perceptrons faded; this contributed to broader skepticism surrounding early AI research.
 \item The first AI winter had \emph{multiple} technical, funding, and expectation-related causes, not one mathematical counterexample.
\end{itemize}
\begin{keyidea}One layer was not expressive enough for every classification problem. The next question: can we combine simple units into a more capable model?\end{keyidea}
\end{frame}

\section{From a Neuron to a Network}
\begin{frame}[t]{A Multilayer Perceptron Combines Learned Features}
\small
\conceptlead{Hidden units can transform the original inputs before the final decision.}
\begin{columns}[T,onlytextwidth]
 \begin{column}{.55\textwidth}
 \lecturevisual[2.64in]{multilayer_network.png}{A fully connected feedforward multilayer network.}
 \end{column}
 \begin{column}{.41\textwidth}
 \begin{itemize}
  \item \textbf{Input layer:} the feature values.
  \item \textbf{Hidden layers:} combinations of earlier signals, followed by nonlinear activations.
  \item \textbf{Output layer:} a class score, class probabilities, or a number.
  \item Stacking transformations can represent boundaries unavailable to one linear unit.
 \end{itemize}
 \end{column}
\end{columns}
\end{frame}

\begin{frame}[t]{How Hidden Units Make XOR Possible}
\small
\conceptlead{A hidden layer can ask simpler questions whose answers solve a harder one.}
\lecturevisual[0.92in]{xor_hidden_features.png}{Two intermediate Boolean patterns: OR and AND.}
One hidden unit detects whether \emph{either} input is on (OR); another detects whether \emph{both} are on (AND). The output combines their answers: \textbf{XOR = OR and not AND}.
\begin{calculation}[A small network implementing XOR with step activations]
\[\begin{aligned}
h_{\rm OR}&=\mathrm{step}(x_1+x_2-0.5),\quad
h_{\rm AND}=\mathrm{step}(x_1+x_2-1.5),\\
\widehat y&=\mathrm{step}(h_{\rm OR}-h_{\rm AND}-0.5).
\end{aligned}\]
$x_1,x_2$: binary inputs; $h_{\rm OR},h_{\rm AND}$: hidden outputs; $\widehat y$: prediction. Check $(1,1)$: both hidden values are 1, so the output is 0. These weights are hand-chosen, not gradient-trained.
\end{calculation}
\end{frame}

\begin{frame}[t]{Inside One Hidden Unit}
\small
\conceptlead{Every hidden unit repeats the same two steps: combine inputs, then activate.}
\begin{enumerate}
 \item Multiply each incoming feature by the weight on its connection.
 \item Add those contributions and a bias.
 \item Apply an activation function to the sum.
 \item Send the resulting value to the next layer.
\end{enumerate}
\begin{calculation}[Forward computation for a hidden unit]
\[z_j=b_j+\sum_k w_{jk}a_k,\qquad a_j=\phi(z_j).\]
$j$: current unit; $k$: previous-layer unit; $a_k$: incoming activation; $w_{jk}$: connection weight; $b_j$: bias; $z_j$: weighted score; $\phi$: nonlinearity; $a_j$: resulting activation. Sum over all incoming units.
\end{calculation}
\smallskip
\textbf{Why nonlinear?} Stacked affine transformations (linear weights plus biases) collapse to one affine transformation. Nonlinear activations change what can be represented.
\end{frame}

\begin{frame}[t]{Activations: Threshold, Sigmoid/Tanh, and ReLU}
\small
\conceptlead{The activation determines how a unit responds to its combined input.}
\lecturevisual[0.81in]{activations.png}{Common activation functions; the middle plot includes sigmoid and tanh.}
\begin{itemize}
 \item \textbf{Step:} a hard 0/1 decision, useful for the original perceptron but unsuitable for ordinary gradient training.
 \item \textbf{Sigmoid/tanh:} smooth nonlinear functions; large positive/negative scores can cause nearly flat responses.
 \item \textbf{ReLU:} leaves positive inputs unchanged and sets negative inputs to zero; a practical default for many hidden layers.
\end{itemize}

\end{frame}

\begin{frame}[t]{Activation Formulas: What Each Curve Computes}
\small
\conceptlead{The same weighted score can produce very different activations.}
\begin{calculation}[Common formulas]
\[\sigma(z)=\frac1{1+e^{-z}},\qquad\tanh(z)=\frac{e^z-e^{-z}}{e^z+e^{-z}},\qquad\mathrm{ReLU}(z)=\max(0,z).\]
$z$: weighted input; $e$: Euler's number; $\sigma$: sigmoid; $\tanh$: hyperbolic tangent; ReLU: rectified linear unit; $\max$: take the larger value. ReLU has a kink at zero; implementations choose a derivative convention there.
\end{calculation}
\begin{block}{Check one input: $z=-1$}
Sigmoid gives about $0.269$; tanh gives about $-0.762$; ReLU gives $0$.
\end{block}
\begin{keyidea}Sigmoid is bounded between 0 and 1; tanh between $-1$ and 1; ReLU is unbounded above. A hidden activation is not automatically a class probability.\end{keyidea}
\end{frame}

\begin{frame}[t]{The Missing Piece: How Do Hidden Layers Learn?}
\small
\conceptlead{The correct output tells us the network is wrong; it does not say which internal weight caused the error.}
\begin{columns}[T,onlytextwidth]
 \begin{column}{.50\textwidth}
 \begin{itemize}
  \item A perceptron's output directly exposes its decision error.
  \item A multilayer network has many hidden weights affecting the same output.
  \item If the final prediction is wrong, should we change an early connection, a late connection, or both?
  \item Hard thresholds offer no useful local slope for gradient-based credit assignment.
 \end{itemize}
 \end{column}
 \begin{column}{.46\textwidth}
 \begin{block}{Historical distinction}
Multilayer networks can \textbf{represent} XOR, even with step activations. But efficient, scalable \textbf{training} requires a way to assign responsibility for error across layers.
 \par\medskip
 \textbf{Representability is not trainability.}
 \end{block}
 \end{column}
\end{columns}
\end{frame}

\section{Backpropagation Makes Hidden Layers Trainable}
\begin{frame}[t]{Backpropagation: Send Error Information Backward}
\small
\conceptlead{Compute how a small change to each weight would affect the current loss.}
\lecturevisual[1.40in]{backprop_flow.png}{Forward prediction and backward gradient propagation.}
\begin{itemize}
 \item \textbf{Forward pass:} compute each activation and the final prediction.
 \item \textbf{Loss:} measure how different the prediction is from the label.
 \item \textbf{Backward pass:} compute the loss gradient with respect to each weight.
 \item \textbf{Optimizer:} use those sensitivities to update the weights.
\end{itemize}
\end{frame}

\begin{frame}[t]{The Chain Rule: Assigning Responsibility}
\small
\conceptlead{A weight affects the loss through a chain of intermediate calculations.}
\begin{itemize}
 \item Picture a weight changing a hidden score, which changes an activation, which changes a prediction, which changes the loss.
 \item A \textbf{derivative} expresses how strongly one value changes as another changes slightly.
 \item Backpropagation multiplies the sensitivities along the chain so each weight gets a meaningful update direction.
\end{itemize}
\begin{calculation}[For a weight $w$ that affects a single downstream activation]
\[\frac{\partial L}{\partial w}
=\underbrace{\frac{\partial L}{\partial a}}_{\text{activation to loss}}
\underbrace{\frac{\partial a}{\partial z}}_{\text{score to activation}}
\underbrace{\frac{\partial z}{\partial w}}_{\text{weight to score}}.\]
$L$: loss; $w$: one weight; $z$: pre-activation score; $a$: activation; $\partial L/\partial w$: local loss sensitivity to that weight, holding other weights fixed. If a value feeds several downstream paths, their gradient contributions are added.
\end{calculation}
\end{frame}

\begin{frame}[t]{Backpropagation and Gradient Descent Have Different Jobs}
\small
\conceptlead{Backpropagation finds the sensitivities; the optimizer decides how to use them.}
\begin{columns}[T,onlytextwidth]
 \begin{column}{.48\textwidth}
 \begin{block}{Backpropagation}
 \textbf{Calculates} how each weight influences the current loss using the chain rule.
 \end{block}
 \end{column}
 \begin{column}{.48\textwidth}
 \begin{block}{Gradient descent / Adam}
 \textbf{Changes} weights in ways expected to lower loss; Adam adapts update sizes using gradient history.
 \end{block}
 \end{column}
\end{columns}
\begin{calculation}[One basic gradient-descent step]
\[w^{(t+1)}=w^{(t)}-\eta\frac{\partial L}{\partial w}.\]
$t$: update index; $w^{(t)},w^{(t+1)}$: old/new weight; $\eta>0$: learning rate; $\partial L/\partial w$: gradient evaluated at the current weights. A batch is a group used for one update; an epoch is one pass through the training examples.
\end{calculation}
\begin{keyidea}Backpropagation is an efficient way to compute gradients for many connected parameters; it is not itself the weight-update policy.\end{keyidea}
\end{frame}

\begin{frame}[t]{1986 and the Return of Multilayer Learning}
\small
\conceptlead{Efficient error propagation made multilayer networks trainable in practice.}
\begin{itemize}
 \item Work on reverse-mode differentiation and error propagation predates 1986.
 \item The influential 1986 work of \textbf{Rumelhart, Hinton, and Williams} helped popularize backpropagation for multilayer connectionist models.
 \item Smooth activations and computed gradients supplied a practical procedure for learning hidden representations.
 \item Better training algorithms did not immediately make deep networks dominant: computational power and useful datasets were still limited.
\end{itemize}
\begin{keyidea}The conceptual breakthrough was an algorithm to \textbf{train} internal layers, not merely a drawing of a deeper network.\end{keyidea}
\end{frame}

\begin{frame}[t]{Why Deep Networks Still Struggled in the 1990s}
\small
\conceptlead{A successful training principle still depends on data, compute, and numerical behavior.}
\begin{itemize}
 \item \textbf{Data:} many tasks lacked large, labeled datasets.
 \item \textbf{Compute:} training many layers with the available hardware was expensive.
 \item \textbf{Vanishing gradients:} saturated sigmoid/tanh units could pass tiny update signals through many layers.
 \item \textbf{Optimization and initialization:} training could proceed slowly or unpredictably.
\end{itemize}
\begin{block}{Historical context}
The late-1980s/early-1990s AI funding downturn had several causes, including disappointment with expert systems. Neural-network research continued; these periods should not be explained by a single architecture or training problem.
\end{block}
\end{frame}

\begin{frame}[t]{The 2010s: Data, Compute, and Better Training}
\small
\conceptlead{Deep learning became practically effective when several enabling conditions came together.}
\begin{columns}[T,onlytextwidth]
 \begin{column}{.53\textwidth}
 \begin{itemize}
  \item \textbf{Larger datasets:} ImageNet, speech corpora, and other large collections supported richer models.
  \item \textbf{Accelerators:} GPUs and distributed computation made larger training runs possible.
  \item \textbf{Improved training:} ReLU, better initialization, normalization, regularization, and more capable optimizers helped stabilize learning.
 \end{itemize}
 \end{column}
 \begin{column}{.43\textwidth}
 \lecturevisual[2.21in]{multilayer_network.png}{The same weighted-layer idea can be repeated at larger scale.}
 \end{column}
\end{columns}
\begin{keyidea}Deep learning was not one new neuron formula: it was old core ideas plus advances in data, hardware, architecture, and optimization.\end{keyidea}
\end{frame}

\section{Practical Python: Breast-Cancer Classification}
\begin{frame}[t]{A Modern Feedforward Network: What Gets Learned?}
\small
\conceptlead{Training adjusts weights and biases; the architecture specifies how they are connected.}
\begin{columns}[T,onlytextwidth]
 \begin{column}{.55\textwidth}
 \lecturevisual[2.60in]{multilayer_network.png}{Fully connected input, hidden, and output layers.}
 \end{column}
 \begin{column}{.41\textwidth}
 \begin{itemize}
  \item \textbf{Parameters:} weights and biases fitted from labeled examples.
  \item \textbf{Hyperparameters:} number of layers, units, activation functions, batch size, epochs, learning rate.
  \item \textbf{Loss:} tells us how far the outputs are from the desired targets.
  \item \textbf{Validation:} helps assess whether improvements extend to unseen examples.
 \end{itemize}
 \end{column}
\end{columns}
\end{frame}

\begin{frame}[t]{Breast-Cancer Example: Reuse the Same Evaluation Protocol}
\small
\conceptlead{Change the model, keeping the data and evaluation rules fixed.}
\begin{itemize}
 \item Use the training, validation, and test arrays prepared before the perceptron example.
 \item Every record is a vector of 30 numerical measurements.
 \item The scaler was fitted only on training rows; validation and test use those same training statistics.
 \item The MLP learns from training rows and monitors validation rows; neither model trains on the final test.
\end{itemize}
\begin{keyidea}This lets us compare model families fairly. The data label $1$ still means benign; increasing a sigmoid score does not mean increasing malignancy risk.\end{keyidea}
\end{frame}

\begin{frame}[fragile,t]{Breast-Cancer Example: Define the Neural Network}
\small
\conceptlead{The architecture stacks two nonlinear hidden layers and one binary output.}
\begin{columns}[T,onlytextwidth]
 \begin{column}{.57\textwidth}
{\footnotesize\begin{verbatim}
import tensorflow as tf
tf.keras.utils.set_random_seed(42)
model = tf.keras.Sequential([
    tf.keras.layers.Input(shape=(30,)),
    tf.keras.layers.Dense(64, activation='relu'),
    tf.keras.layers.Dense(32, activation='relu'),
    tf.keras.layers.Dense(1, activation='sigmoid')
])
\end{verbatim}}
 \end{column}
 \begin{column}{.40\textwidth}
 \begin{itemize}
  \item \texttt{Input(30)}: 30 measurements.
  \item \texttt{Dense(64)}: 64 learned feature combinations.
  \item \texttt{Dense(32)}: 32 more transformations.
  \item \texttt{Dense(1)} with sigmoid: one score between 0 and 1.
 \end{itemize}
 \end{column}
\end{columns}
\end{frame}

\begin{frame}[fragile,t]{Breast-Cancer Example: Compile the Model}
\small
\conceptlead{Compilation specifies what error means and how the weights will change.}
{\footnotesize\begin{verbatim}
model.compile(
    optimizer='adam',
    loss='binary_crossentropy',
    metrics=['accuracy']
)
\end{verbatim}}
\begin{itemize}
 \item \textbf{Loss:} binary cross-entropy penalizes confident probability mistakes.
 \item \textbf{Optimizer:} Adam uses backpropagated gradients to update trainable parameters.
 \item \textbf{Metric:} accuracy counts correct labels at the evaluation threshold.
\end{itemize}
\begin{calculation}[Loss for one binary-labeled observation]
\[\ell(p,y)=-\bigl[y\ln p+(1-y)\ln(1-p)\bigr].\]
$y\in\{0,1\}$: true label; $p\in(0,1)$: predicted probability of class 1; $\ell$: one-example loss; $\ln$: natural logarithm. For $y=1$, loss is $-\ln p$; for $y=0$, it is $-\ln(1-p)$. $L$ usually averages such losses over a batch.
\end{calculation}
\end{frame}

\begin{frame}[fragile,t]{Breast-Cancer Example: Fit for 30 Epochs}
\small
\conceptlead{An epoch is one pass through the training examples; batches make several updates per epoch.}
{\footnotesize\begin{verbatim}
history = model.fit(
    X_train, y_train,
    epochs=30,
    batch_size=16,
    validation_data=(X_val, y_val),
    verbose=1
)
\end{verbatim}}
\begin{itemize}
 \item \texttt{batch\_size=16}: update weights after processing groups of sixteen examples.
 \item \texttt{validation\_data}: monitor the separate validation rows without training on them.
 \item \texttt{history}: stores epoch-by-epoch loss and accuracy for plotting.
\end{itemize}
\begin{keyidea}Freeze settings chosen using validation, then report \texttt{model.evaluate(X\_test, y\_test)} once.\end{keyidea}
\end{frame}

\begin{frame}[t]{Reading Training Histories}
\small
\conceptlead{Look for generalization, not only rising training accuracy.}
\lecturevisual[1.08in]{training_curves_illustrative.png}{Illustrative training and validation curves (not a measured Keras run).}
\begin{itemize}
 \item Accuracy should improve as useful relationships are learned, but fluctuations are normal.
 \item If training loss falls while validation loss rises, the model may be overfitting.
 \item A high validation accuracy in a small sample is not a guarantee of future reliability.
\end{itemize}
\begin{calculation}[Retrieving the arrays returned by Keras]
\texttt{history.history['loss']} and \texttt{history.history['val\_loss']} are loss curves; the \texttt{accuracy} and \texttt{val\_accuracy} entries are classification curves.
\end{calculation}
\end{frame}

\begin{frame}[t]{Compare the Supplied Training History Carefully}
\small
\conceptlead{A supplied result illustrates a run; it is not a guaranteed outcome of our revised split.}
\lecturevisual[1.62in]{source_breast_cancer_results.png}{Training and validation history supplied by the instructor.}
\begin{itemize}
 \item Both validation curves improve here; compare this with the illustrative overfitting curves on the previous slide.
 \item Split choices, preprocessing, random seeds, and training details affect the numbers.
\end{itemize}
\end{frame}

\section{What Networks Can Approximate}
\begin{frame}[t]{Continuous Outputs and Universal Approximation}
\small
\conceptlead{A neural network is not limited to classification: it can predict continuous quantities.}
\begin{itemize}
 \item Replacing the final sigmoid with a linear output lets a network predict a numerical value.
 \item With a suitable nonlinear activation and enough hidden units, certain neural networks can approximate any continuous target on a bounded input region.
 \item This is an \textbf{existence result}, not a promise that a particular dataset or optimizer will find the desired approximation.
 \item Depth can provide more efficient representations for some structured functions; it is not automatically superior on every task.
\end{itemize}
\begin{calculation}[The theorem's practical reading]
For a continuous target $g$ on a closed, bounded region $K$, suitable networks can be constructed such that
\[\max_{x\in K}|f_{\theta}(x)-g(x)|<\varepsilon\qquad\text{for any chosen }\varepsilon>0.\]
$K$: compact (closed, bounded) input set; $x$: input; $g$: continuous target; $f_\theta$: network output; $\theta$: all weights and biases; $\varepsilon>0$: chosen error tolerance. $\max$ measures the worst absolute error throughout $K$. The result does not guarantee a feasible size or a successful training procedure.
\end{calculation}
\end{frame}

\begin{frame}[fragile,t]{Original Python Example: Learning a Sine Wave}
\small
\conceptlead{Use the same feedforward architecture for a numerical target instead of a class label.}
{\scriptsize\begin{verbatim}
import numpy as np
import tensorflow as tf

X = np.linspace(0, 2 * np.pi, 1000).reshape(-1, 1)
y = np.sin(X)
model = tf.keras.Sequential([
    tf.keras.layers.Input(shape=(1,)),
    tf.keras.layers.Dense(64, activation='tanh'),
    tf.keras.layers.Dense(64, activation='tanh'),
    tf.keras.layers.Dense(1)
])
model.compile(optimizer='adam', loss='mse')
model.fit(X, y, epochs=100, batch_size=32, verbose=0)
\end{verbatim}}
\begin{keyidea}The last layer has \textbf{no sigmoid}, so it can produce any real number. Mean squared error measures numerical prediction differences.\end{keyidea}
\end{frame}

\begin{frame}[t]{A Network Can Fit a Smooth Function}
\small
\conceptlead{The network learns to combine smooth nonlinear transformations into a smooth prediction.}
\begin{columns}[T,onlytextwidth]
 \begin{column}{.57\textwidth}
 \lecturevisual[2.53in]{sine_approx_illustrative.png}{Illustrative sine approximation, not the output of the original Keras run.}
 \end{column}
 \begin{column}{.39\textwidth}
 \begin{itemize}
  \item The input is one position $x$; the target is $\sin(x)$. The supplied curve is illustrative, not a rerun.
  \item Two hidden layers with \texttt{tanh} allow a curved numerical relationship.
  \item Training MSE asks how well the approximation matches observed targets.
  \item Evaluation on new inputs is needed to assess interpolation and generalization.
 \end{itemize}
 \end{column}
\end{columns}
\end{frame}

\begin{frame}[t]{MNIST: Recognizing Handwritten Digits}
\small
\conceptlead{The same model family can classify images by treating their pixel values as features.}
\begin{columns}[T,onlytextwidth]
 \begin{column}{.52\textwidth}
 \lecturevisual[2.52in]{source_mnist_grid.png}{Examples from the instructor's original MNIST slide.}
 \end{column}
 \begin{column}{.44\textwidth}
 \begin{itemize}
  \item Each grayscale image has $28\times28=784$ pixels.
  \item Flatten the image to a vector of 784 input values.
  \item Normalize intensities from 0--255 into the 0--1 range.
  \item Predict one of ten classes: the digits 0 through 9.
 \end{itemize}
 \end{column}
\end{columns}
\end{frame}

\begin{frame}[fragile,t]{Original Python Example: A Dense MNIST Classifier}
\small
\conceptlead{Replace the binary sigmoid with a ten-way softmax output.}
{\scriptsize\begin{verbatim}
(X_train, y_train), (X_test, y_test) = (
    tf.keras.datasets.mnist.load_data())
X_train = X_train.reshape(-1, 784).astype('float32') / 255.0
X_test  = X_test.reshape(-1, 784).astype('float32') / 255.0
model = tf.keras.Sequential([
    tf.keras.layers.Input(shape=(784,)),
    tf.keras.layers.Dense(128, activation='relu'),
    tf.keras.layers.Dense(64, activation='relu'),
    tf.keras.layers.Dense(10, activation='softmax')
])
\end{verbatim}}
\begin{calculation}[Why softmax at the output?]
\[p_k=\frac{e^{z_k}}{\sum_{j=0}^{9}e^{z_j}},\qquad \sum_{k=0}^{9}p_k=1.\]
$z_k$: raw output score (logit) for digit $k$; $e$: Euler's number; $j,k\in\{0,\ldots,9\}$: class indices; $p_k$: estimated class probability. Choose the largest $p_k$; normalization alone does not guarantee calibration.
\end{calculation}
\end{frame}

\begin{frame}[fragile,t]{MNIST: Train With Integer Labels and Test Once}
\small
\conceptlead{Keep a validation subset of the 60,000 training images; preserve the official test set.}
{\footnotesize
\begin{verbatim}
X_train, X_val, y_train, y_val = train_test_split(
    X_train, y_train, test_size=10000,
    stratify=y_train, random_state=42)
model.compile(optimizer='adam',
    loss='sparse_categorical_crossentropy', metrics=['accuracy'])
model.fit(X_train, y_train, epochs=3, batch_size=128,
    validation_data=(X_val, y_val), verbose=0)
print("MNIST test:", model.evaluate(X_test, y_test, verbose=0))
\end{verbatim}}
\begin{itemize}
 \item \texttt{sparse\_categorical\_crossentropy} expects integer labels $0$ through $9$, not one-hot vectors.
 \item Three epochs is a preset demonstration budget; use validation if you change it.
\end{itemize}
\begin{keyidea}The supplied confusion matrix is a separate earlier result. A newly trained model must be evaluated on its own predictions.\end{keyidea}
\end{frame}

\begin{frame}[t]{MNIST: Evaluate Confusion, Not Only Accuracy}
\small
\conceptlead{A high overall score can still hide systematic confusion between particular digits.}
\begin{columns}[T,onlytextwidth]
 \begin{column}{.53\textwidth}
 \lecturevisual[2.08in]{source_mnist_confusion.png}{Confusion matrix displayed in the instructor's original MNIST example.}
 \end{column}
 \begin{column}{.43\textwidth}
 \begin{itemize}
  \item \textbf{Diagonal cells:} images assigned the correct digit.
  \item \textbf{Off-diagonal cells:} confusion between two digits.
  \item This supplied chart is not a result of our new run. Axis direction and counts must be labeled before reporting which digits are confused.
  \item Flattening retains pixel positions as vector indices. Dense layers do not build in local neighborhoods or shared filters; CNNs do.
 \end{itemize}
 \end{column}
\end{columns}
\end{frame}

\section{Limitations and Model Choices}
\begin{frame}[t]{Downsides of Neural Networks}
\small
\conceptlead{More expressive models come with costs and failure modes.}
\begin{columns}[T,onlytextwidth]
 \begin{column}{.49\textwidth}
 \begin{itemize}
  \item Large labeled datasets may be needed to fit many parameters reliably.
  \item Deep internal representations are harder to interpret than short decision rules.
  \item Larger networks may demand substantial computation and memory.
 \end{itemize}
 \end{column}
 \begin{column}{.48\textwidth}
 \begin{itemize}
  \item Overfitting can still occur, especially with small or noisy samples.
  \item Prediction confidence is not automatically calibrated uncertainty.
  \item Training depends on architecture, initialization, learning rate, and other choices.
 \end{itemize}
 \end{column}
\end{columns}
\lecturevisual[0.96in]{neural_overfitting.png}{Illustrative stable versus noise-sensitive fits.}
\end{frame}

\begin{frame}[t]{Hyperparameters: What We Choose Before Training}
\small
\conceptlead{Learn weights from data; use validation evidence to choose the training design.}
\begin{center}
\renewcommand{\arraystretch}{1.13}
\begin{tabular}{>{\raggedright\arraybackslash}p{.26\linewidth}>{\raggedright\arraybackslash}p{.34\linewidth}>{\raggedright\arraybackslash}p{.30\linewidth}}
\toprule
\textbf{Choice} & \textbf{What it controls} & \textbf{Common examples}\\
\midrule
Layers, units & Network depth and width & 2 layers; 64 or 128 units\\
Activation & How a neuron transforms input & ReLU, tanh, sigmoid\\
Learning rate / optimizer & How fast and how weights change & Adam, SGD\\
Batch size / epochs & Updates per pass; passes through data & 16, 32 / 30, 100\\
Regularization & Controls overly complex fits & Dropout, weight decay\\
Stopping & When training should end & Early stopping\\
\bottomrule
\end{tabular}
\end{center}
\begin{keyidea}Use training data to \textbf{fit} weights, validation data to \textbf{choose} hyperparameters, and a final untouched test set to \textbf{evaluate} the selected model.\end{keyidea}
\end{frame}

\begin{frame}[t]{From One Perceptron to Deep Learning}
\small
\begin{enumerate}
 \item A perceptron learns a linear boundary by changing feature weights after mistakes.
 \item Hidden layers make more complex functions possible, including XOR.
 \item Nonlinear activation functions provide useful intermediate representations and gradients.
 \item Backpropagation calculates weight sensitivities; the optimizer updates the parameters.
 \item Data, accelerators, and improved training methods made deep learning practical at larger scale.
 \item The same basic framework supports binary classification, continuous regression, and multiclass image recognition.
\end{enumerate}
\begin{keyidea}Representability asks \emph{what the network could express}; optimization asks \emph{what training can find}; generalization asks \emph{whether it works on new cases}.\end{keyidea}
\smallskip
\centering\scriptsize Next: architectures that exploit structure in images and sequences (CNNs and RNNs); generative networks are another extension.
\end{frame}
\end{document}
